% !TEX root = ../main.tex

\chapter{Introduction}\label{chapter:introduction}

Urban transportation networks are at once deeply familiar and poorly understood: millions of daily journeys produce emergent traffic patterns that resist simple analytical description~\cite{horni2016matsim}. When a city government considers closing a major arterial, reducing road capacity to encourage modal shift, or investing in new infrastructure, it faces a fundamental question: what will actually happen to traffic? Agent-based models (ABMs) have become the standard methodology for answering such questions, simulating the decisions of individual travellers to predict how policy interventions reshape flow across the network~\cite{bonabeau2002agent, railsback2019agent}.

The computational demands of these simulations, however, create a real bottleneck in the policy evaluation process. Simulating a single scenario for a large metropolitan area like Paris can require eight or more hours on modern hardware---severely limiting the number of alternatives that planners can feasibly explore within project timelines~\cite{natterer2025ml}.

Machine learning surrogates can address this bottleneck. By training neural networks to approximate simulation outputs, scenario evaluation times can be reduced from hours to seconds~\cite{willard2022integrating}. Recent work by Natterer~et~al.~\cite{natterer2025ml} demonstrated that Graph Neural Networks (GNNs) are particularly well-suited for this task: they naturally capture the graph structure of road networks while learning complex relationships between policy interventions and traffic flow changes.

Yet accuracy alone does not suffice for decision support. When a surrogate model encounters scenarios that differ substantially from its training distribution, its predictions may become unreliable---and without uncertainty quantification (UQ), users have no principled way to identify these cases. Without such estimates, planners cannot distinguish reliable predictions from unreliable ones.

\begin{figure}[htbp]
\centering
\includegraphics[width=0.95\textwidth]{figures/fig_network_intro.pdf}
\caption[Paris road network as directed graph for the GNN surrogate]{%
    Directed line-graph representation of the Paris road network used by the GNN surrogate.
    Each node represents one road segment with five input features (listed right);
    coral-shaded nodes carry a non-zero \texttt{CAPACITY\_REDUCTION} due to the policy
    intervention.  The per-node prediction target is
    $\Delta v = v_{\mathrm{policy}} - v_{\mathrm{baseline}}$ (veh/h).
    The Paris network contains 31{,}635 nodes.
    Source:~\textcite{natterer2025ml}.}
\label{fig:network_schematic}
\end{figure}

This thesis investigates that gap. We study uncertainty quantification methods for GNN-based traffic surrogates, conducting a systematic comparison of Monte Carlo (MC) Dropout against ensemble-based approaches and conformal prediction. The goal is to establish which UQ methods are most effective for graph-structured traffic prediction, and to provide practical guidance for deploying uncertainty-aware surrogates in urban planning. All experiments use 1,000 of 10,000 available MATSim scenarios (10\% subset).

\section{Problem Statement}

GNN surrogates for traffic prediction must generalise to policy scenarios that may differ substantially from their training distribution. In classical supervised learning, this distribution shift manifests as degraded prediction accuracy. In decision support applications, however, the consequences are more severe---planners may unknowingly base infrastructure decisions on unreliable predictions.

The research gap we address comprises three interconnected components:

\begin{enumerate}
    \item \textbf{Epistemic uncertainty in GNN surrogates:} While uncertainty quantification has been studied in computer vision and natural language processing, its application to GNN-based traffic prediction has received little attention. The graph structure of road networks introduces unique challenges, as uncertainty may propagate through connected road segments.
    
    \item \textbf{Method comparison:} MC Dropout and deep ensembles represent the two dominant paradigms for uncertainty quantification in deep learning~\cite{gal2016dropout, lakshminarayanan2017simple}. However, their comparative effectiveness for graph-structured traffic data has not been systematically evaluated.
    
    \item \textbf{Practical deployment guidance:} Urban planners require concrete recommendations about which UQ method to adopt, how to configure its parameters, and how to interpret the resulting uncertainty estimates in decision-making contexts.
\end{enumerate}

\section{Research Questions}

This thesis addresses four closely related research questions:

\begin{description}
    \item[RQ1:] How effectively does MC Dropout capture epistemic uncertainty in GNN-based traffic surrogates, as measured by the correlation between predicted uncertainty and actual prediction error?
    
    \item[RQ2:] How does MC Dropout compare to ensemble-based uncertainty quantification---both variance across training runs and multi-model ensembles---in terms of uncertainty quality and computational efficiency?
    
    \item[RQ3:] Does combining uncertainty estimates from architecturally identical models provide benefits over single-model UQ methods, and what role does model diversity play in ensemble effectiveness?
    
    \item[RQ4:] Can distribution-free methods---specifically conformal prediction and post-hoc calibration---transform the raw uncertainty estimates into trustworthy prediction intervals with formal coverage guarantees?
\end{description}

\section{Contributions}

The contributions of this thesis are as follows:

\begin{enumerate}
    \item \textbf{Systematic UQ evaluation:} We conduct the first comprehensive comparison of uncertainty quantification methods for GNN-based traffic surrogates, evaluating MC Dropout, ensemble variance, and multi-model ensembles on a large-scale Paris road network dataset with 3.16 million node-level predictions across 100 test scenarios.
    
    \item \textbf{Empirical findings on method effectiveness:} For Trial~8 (best model), MC Dropout achieves Spearman~$\rho = 0.4820$ between predicted uncertainty and absolute error (100 test graphs, 3,163,500 predictions). In Experiment~A (5 independently seeded MC~Dropout inference runs), combining MC~Dropout and ensemble variance yields $\rho = 0.4909$, marginally exceeding standalone MC~Dropout.  Multi-model ensemble disagreement (Experiment~B, 5 models) achieves $\rho = 0.4333$---useful but weaker than MC~Dropout.
    
    \item \textbf{Ensemble evaluation:} Ensemble experiments confirm that MC~Dropout ($\rho = 0.4908$) outperforms seed-based ensemble variance ($\rho = 0.4370$) by 12.3\%, while multi-model disagreement ($\rho = 0.4333$) provides complementary but weaker uncertainty ranking. A PyTorch Geometric API version mismatch originally corrupted these results; the corrected evaluation (Section~\ref{sec:ensemble_caveat}) resolves the issue and enables a reliable comparison.
    
    \item \textbf{Deployment recommendations:} We provide concrete guidance for practitioners implementing uncertainty quantification in traffic GNN systems, including parameter choices ($S=30$ MC samples), model selection strategies, and frameworks for interpreting uncertainty estimates.
    
    \item \textbf{Post-hoc calibration and coverage guarantees:} We demonstrate that split conformal prediction achieves near-exact marginal coverage (90.02\% and 95.01\% at the respective nominal levels) on 1.58M-node evaluation sets, and that post-hoc temperature scaling reduces Expected Calibration Error by 82\%. The combination of MC Dropout (for uncertainty ranking), temperature scaling (for improved calibration), and conformal prediction (for formal coverage guarantees) forms a practical, layered UQ framework requiring no model retraining.
\end{enumerate}

\section{Thesis Structure}

The remainder of this thesis is organized as follows:

\textbf{Chapter~\ref{chapter:background}: Background and Related Work} establishes the theoretical foundations of this research. Section~\ref{sec:abm} introduces agent-based transport models with emphasis on the MATSim framework. Section~\ref{sec:gnn} presents the fundamentals of graph neural networks and their application to traffic prediction. Section~\ref{sec:uq} reviews uncertainty quantification methods in deep learning, including Bayesian approaches and ensemble methods.

\textbf{Chapter~\ref{chapter:methodology}: Methodology} describes the technical approach adopted in this thesis. We present the PointNetTransfGAT architecture, detail our implementations of MC Dropout and ensemble-based UQ methods, and define the evaluation metrics used to assess uncertainty quality.

\textbf{Chapter~\ref{chapter:experiments}: Experimental Setup} presents the Paris road network dataset, training protocol, and experimental design. We describe two complementary experiments: Experiment~A comparing MC Dropout against ensemble variance on a single model, and Experiment~B evaluating multi-model ensemble approaches.

\textbf{Chapter~\ref{chapter:results}: Results} reports the findings from both experiments, including quantitative comparisons of UQ methods, conformal prediction coverage analysis, post-hoc calibration results, and analysis of the generalisation gap observed between validation and test performance.

\textbf{Chapter~\ref{chapter:discussion}: Discussion} interprets results in the context of prior work, addresses limitations of our study, and discusses implications for deploying uncertainty-aware traffic prediction systems.

\textbf{Chapter~\ref{chapter:conclusion}: Conclusion} summarizes the key findings, revisits the research questions, and outlines directions for future work.

\begin{figure}[htbp]
\centering
\includegraphics[width=\textwidth]{figures/fig11_thesis_workflow.pdf}
\caption[Overall thesis pipeline: from MATSim scenarios to policy decisions]{Overall thesis pipeline. MATSim agent-based simulations (1,000 of 10,000 available scenarios, 10\% subset) are converted to line-graph representations and used to train eight GNN trials. The best model (T8, $R^2 = 0.5957$) is equipped with MC Dropout uncertainty ($\rho = 0.4820$) and conformal prediction intervals (90\%/95\% coverage guaranteed), feeding into an uncertainty-guided policy decision framework.}
\label{fig:thesis_workflow}
\end{figure}
